\section*{历史注记}
\phantomsection
\addcontentsline{toc}{section}{历史注记}

（括号中的数字指本注记末尾的参考文献。）

如我们在第 II 章的历史注记中所说，度量空间的概念是 M. Fréchet 于 1906 年引入的，几年后由 F. Hausdorff 在其``Mengenlehre''（《集论》）中加以发展。1920 年以后它获得了巨大的重要性，部分地是由于 S. Banach 及其学派关于赋范空间及其在泛函分析中的应用的基础性工作，也是由于绝对值概念在数论与代数几何中的重要性（在那里，关于一个绝对值的完备化过程特别地显示出是一种有力的工具）。

1920-1930 这十年间出现了一整系列关于度量空间性质的研究。这些由莫斯科学派进行的研究，特别着眼于寻求给定的拓扑可度量化的必要充分条件。这一思想运动凸显了正规空间概念的意义，该概念由 Tietze 于 1923 年定义，但其重要作用只是作为 Urysohn 关于连续实值函数扩张的工作{[}6{]}的结果才被认识到。除去单实变量函数这一平凡情形外，把定义在闭集上的连续实值函数扩张到整个空间的问题，最先（就平面的情形）由 H. Lebesgue{[}3{]}考虑；在 Urysohn 的决定性结果之前，它已由 H. Tietze{[}4{]}对度量空间解决。把这问题推广到取值于任意拓扑空间的函数，近年来在代数拓扑中获得了相当的重要性。最近的工作还表明，在这类问题中正规空间的概念不易处理，因为它容许太多``病态''的可能；它常常必须用更受限制的仿紧性概念来代替，后者是 J. Dieudonné 于 1944 年引入的{[}9{]}。在这一理论中，最值得注意的结果是 A. H. Stone 的定理{[}10{]}：每个可度量化空间都是仿紧的 (*)。

我们已经指出过（第 IV 章的历史注记），十九世纪末与二十世纪初（E. Borel、Baire、Lebesgue、Osgood、W. H. Young）关于 \(R^{n}\) 中点集的分类，以及关于通过迭代取极限过程（关于函数列）从连续函数得到的实值函数的分类与刻画的重要工作。人们很快认识到，度量空间——其 1910 年以后的发展主要是俄国与波兰学派的工作——构成了这类研究的自然领域。除其他方面外，度量空间的研究清楚地表明了贫集概念以及完备度量空间中稠密开集的可数交定理（ \(§\ 5,\ no.\ 3,\ Theorem\ 1\) ）在现代分析中所起的基本作用，该定理最先（相互独立地）由 Osgood{[}1{]}对实直线证明，并由 Baire{[}2{]}对空间 \(R^{n}\) 证明。

另一方面，Souslin 于 1917 年{[}5{]}在纠正 Lebesgue 的一个错误时证明了波莱尔集的连续像未必是波莱尔集；这引导他定义并研究更大的一类集，此后称为``解析集''或``苏斯林集''。Souslin 早逝之后，这项工作特别由 N. Lusin（其思想曾启发 Souslin 的工作）与波兰数学家们（见{[}7{]}与{[}8{]}）继续推进。这些集如今的重要性在于它们对积分理论（在那里，由于它们的特殊性质，可以进行对任意可测集而言不可能的构造）以及现代位势论的应用；在位势论中，关于苏斯林集可容性的基本定理（最近由 G. Choquet{[}11{]}证明）已显示出在种种应用中的丰富内容。

